\documentclass{aastex701}
\journalinfo{}

\PassOptionsToPackage{hypertexnames=false}{hyperref}
\PassOptionsToPackage{noabbrev}{cleveref}
\PassOptionsToPackage{noabbrev}{cleveref}

\usepackage{amssymb}
\usepackage{textcomp}
\usepackage{fix-cm}
\usepackage{type1cm}
\usepackage{lmodern}
% Remove the specific nameref warning by temporarily restoring the kernel \label
% before hyperref/nameref is initialized, then restore the class's \label.
\makeatletter
\let\AAS@label\label
\let\label\ltx@label
\makeatother

% Some engines/classes do not define \euro; keep a safe alias.
\providecommand{\euro}{\texteuro}

% Silence the specific nameref warning about \label redefinition.
\usepackage{silence}
\WarningFilter{nameref}{The definition of \\label has changed!}

\begin{document}

\title{Economic Observability (Non-Normative)}

\author{Egor Merkushev}
\email{gorkaedeep@gmail.com}
\affiliation{Independent Researcher}

\begin{abstract}
We introduce a non-normative formulation for observing and comparing the
economic footprint of AI-mediated software construction across specification
graph nodes and time windows. The goal is diagnostic: to expose inference,
verification, and ownership proxies as observable quantities under a chosen
pricing surface, enabling longitudinal and comparative analysis rather than
optimization prescriptions.
\end{abstract}

\keywords{software engineering --- AI-assisted programming --- software metrics --- observability --- cost modeling}

\section{Economic Observability}

\subsection{Pricing Surface and Footprints}

Let each specification node be denoted by $u \in V$, and let $t$ denote an
iteration window (e.g., a generation run, a PR cycle, or a bounded time slice).

We assume a configurable \emph{pricing surface} for the active toolchain:

\begin{itemize}
\item $\mathcal{M}$ --- the set of AI models available in the product.
\item $\mathcal{J}$ --- the set of external tools (search, build, test, OCR, etc.).
\item $p_m$ --- price per token (or billable token unit) for model $m \in \mathcal{M}$.
\item $q_m$ --- fixed per-call price for model $m$ (optional; $q_m=0$ if none).
\item $r_j$ --- price per invocation for tool $j \in \mathcal{J}$ (optional).
\end{itemize}

For each node $u$ and window $t$, we observe the footprint tuple
\[
\left( T_m(u,t),\, N_m(u,t) \right)_{m \in \mathcal{M}},\quad
\left( N_{\mathrm{tool},j}(u,t) \right)_{j \in \mathcal{J}},
\]
where $T_m(u,t)$ is the (billable) token volume under model $m$,
$N_m(u,t)$ is the number of model calls, and $N_{\mathrm{tool},j}(u,t)$ is
the number of invocations of external tool $j$.

\subsection{Node Inference Cost}

We define the \emph{node inference cost} as the pricing-functional applied to
the observed footprint:

\begin{equation}
C_{\mathrm{inf}}(u,t) \equiv
\sum_{m \in \mathcal{M}}
\Bigl(p_m\,T_m(u,t) + q_m\,N_m(u,t)\Bigr)
\;+\;
\sum_{j \in \mathcal{J}} r_j\,N_{\mathrm{tool},j}(u,t).
\end{equation}

This value is directly displayable on the map as currency (\$, \euro, etc.)
under a fixed pricing surface.

\subsection{Graph (Subgraph) Inference Cost}

For any subgraph or selection $S \subseteq V$ (e.g., a feature branch, epic,
release slice), we define:

\begin{equation}
C_{\mathrm{inf}}(S,t) \equiv \sum_{u \in S} C_{\mathrm{inf}}(u,t).
\end{equation}

\subsection{Differential Economics}

Economic observability is intended to be interpreted comparatively. We define
the discrete delta and drift-velocity operators for inference cost:

\begin{equation}
\Delta C_{\mathrm{inf}}(u; t_k) \equiv
C_{\mathrm{inf}}(u,t_k) - C_{\mathrm{inf}}(u,t_{k-1}),
\end{equation}

\begin{equation}
\mathrm{DV}_{C}(u,t) \equiv \frac{d\,C_{\mathrm{inf}}(u,t)}{dt},
\qquad
\mathrm{DV}_{C}(u,t_k) \approx \frac{C_{\mathrm{inf}}(u,t_k)-C_{\mathrm{inf}}(u,t_{k-1})}{t_k-t_{k-1}}.
\end{equation}

Analogous operators apply to any node- or graph-level metric.

\subsection{Verification Cost}

We define a separate (configurable) verification cost functional:

\begin{equation}
C_{\mathrm{verify}}(u,t) \equiv
\sum_{\ell \in \mathcal{L}} s_\ell\,N_{\mathrm{verify},\ell}(u,t),
\end{equation}

where $\mathcal{L}$ is the set of verifiers (CI workflows, device farm runs,
property test suites, performance regressions, etc.),
$N_{\mathrm{verify},\ell}(u,t)$ is the number of verification runs of type $\ell$,
and $s_\ell$ is the unit price for $\ell$ (or an internal cost proxy).

\subsection{Total Build Cost}

We define the \emph{build cost} for a node or selection as:

\begin{equation}
C_{\mathrm{build}}(u,t) \equiv C_{\mathrm{inf}}(u,t) + C_{\mathrm{verify}}(u,t),
\qquad
C_{\mathrm{build}}(S,t) \equiv \sum_{u \in S} C_{\mathrm{build}}(u,t).
\end{equation}

$C_{\mathrm{build}}$ is the direct economic footprint of producing a concrete
artifact state from intent under the current toolchain configuration.

\subsection{Change-Cost Proxy (Graph-Based)}

We introduce a non-normative \emph{change-cost proxy} to quantify how costly
future modifications may become due to graph coupling and instability.
Let $\mathrm{WCR}_k(u)$ denote the weighted change radius within hop distance $k$,
and let $\mathrm{SVC}(u)$ and $\mathrm{CSI}(u)$ denote intent verifiability and
checkpoint stability, respectively (defined elsewhere).

\begin{equation}
C_{\mathrm{change}}(u) \equiv
\kappa\,\mathrm{WCR}_k(u)\,
\frac{1}{\mathrm{SVC}(u)+\varepsilon}\,
\frac{1}{\mathrm{CSI}(u)+\varepsilon},
\end{equation}

where $\kappa$ is a project-calibration coefficient and $\varepsilon>0$ stabilizes
the expression. This quantity is diagnostic: its meaning lies in relative
comparisons and longitudinal drift.

\subsection{Total Ownership Proxy}

Finally, we define a composite \emph{ownership proxy} over a selection $S$:

\begin{equation}
C_{\mathrm{own}}(S) \equiv
C_{\mathrm{build}}(S) + \lambda \sum_{u \in S} C_{\mathrm{change}}(u),
\end{equation}

where $\lambda$ encodes the chosen support horizon (how much weight is assigned
to future change cost). $C_{\mathrm{own}}$ is not an estimate guarantee; it is an
economic lens for comparing structural alternatives and observing drift under
AI-mediated construction.

\subsection{Interpretation Note}

All economic quantities are non-normative observables. Absolute values depend on
pricing, tooling, and policy configuration. The recommended mode of use is
comparative: deltas ($\Delta$), drift velocities ($\mathrm{DV}$), distribution
shifts across the graph, and cost bands under controlled configurations.

\subsection{Map View Encoding (UI Mapping)}

We treat the map as a \emph{projection} of observability signals onto visual channels.
Let $x(u)\in\mathbb{R}^2$ denote the fixed layout coordinate of node $u$ in the map view.
Under a selected scope $S\subseteq V$ and time window $t$, the UI encodes:

\begin{itemize}
\item \textbf{Fill color (economic heat):}
  node fill is mapped to instantaneous inference cost
  \begin{equation}
  \mathrm{Color}(u) \Leftarrow C_{\mathrm{inf}}(u,t),
  \end{equation}
  typically via a monotone colormap with optional banding.

\item \textbf{Node radius (growth / pressure):}
  node size reflects a chosen differential economic signal, by default
  \begin{equation}
  \mathrm{Radius}(u) \Leftarrow \left|\Delta C_{\mathrm{inf}}(u; t_k)\right|
  \quad\text{or}\quad
  \mathrm{Radius}(u) \Leftarrow \left|\mathrm{DV}_{C}(u,t)\right|.
  \end{equation}

\item \textbf{Halo / influence ring (blast radius):}
  an outer ring visualizes change impact at hop distance $k$,
  \begin{equation}
  \mathrm{Halo}(u) \Leftarrow \mathrm{WCR}_k(u),
  \end{equation}
  where the ring radius or intensity increases with weighted reachability.

\item \textbf{Border style (confidence / stability):}
  border opacity or dash pattern encodes verification and stability,
  \begin{equation}
  \mathrm{Border}(u) \Leftarrow \bigl(\mathrm{SVC}(u),\,\mathrm{CSI}(u,t)\bigr),
  \end{equation}
  where low $\mathrm{SVC}$ (weak verifiability) and low $\mathrm{CSI}$ (unstable checkpoints)
  are rendered with reduced opacity and/or higher visual ``noise''.

\item \textbf{Edge emphasis (coupling cost flow):}
  when edges are shown, thickness or glow corresponds to the marginal
  cost transfer associated with coupling, e.g.
  \begin{equation}
  \mathrm{EdgeWeight}(u,v) \Leftarrow
  \Delta C_{\mathrm{inf}}(u; t_k)\,\mathbf{1}_{(u\rightarrow v)}
  \quad\text{or}\quad
  w(u,v),
  \end{equation}
  where $w(u,v)$ is an application-defined coupling weight (dependency strength,
  co-change frequency, or interface-touch severity).

\item \textbf{Scope totals (header KPI, non-normative):}
  the current selection $S$ reports aggregated economics
  \begin{equation}
  \mathrm{Header}(S) \Leftarrow
  \Bigl(C_{\mathrm{inf}}(S,t),\, C_{\mathrm{build}}(S,t),\, C_{\mathrm{own}}(S)\Bigr),
  \end{equation}
  always accompanied by the active pricing surface and time window.
\end{itemize}

The mapping is explicitly \emph{lens-based}: the user may choose which economic signal drives each channel (color, radius, halo, border), while the system
preserves comparability by fixing the pricing surface and time window for a
given analysis session.

\end{document}
